\noindent\textbf{Abstract}\par\smallskip
\noindent Force-plate reports reduce each countermovement jump to a few numbers. Using \nTrials{} routine jumps from \nAth{} male football players, mostly adolescents, we located where along take-off and landing force curves differ. Function-on-scalar regression and athlete-level permutation tests showed that better jumpers push harder in the final third of take-off; elsewhere jump height, age and body mass explained little. Jumps of equal height varied almost as much as the whole sample. Age differences reflected size and performance. Landing profiles formed a continuum but persisted across sessions, so curve features add information that conventional metrics miss.

\medskip
\noindent\textbf{Keywords:} Countermovement jump; Force plates; Functional data analysis; Statistical parametric mapping; Athlete monitoring; Landing; Youth football; Sports analytics
